% Figure: the inverse relation between an antenna's directivity and its
% beamwidth. For a pencil beam the directivity is roughly the sphere's
% solid angle divided by the beam's, D0 ~ 4 pi / (theta_E theta_H) with the
% beamwidths in radians, or the working shorthand D0[dBi] ~ 10 log10(41000/
% (theta_E * theta_H)) with the beamwidths in degrees. We plot directivity
% in dBi against a single symmetric beamwidth theta (degrees), the curve
% that lets an engineer read gain off a measured beamwidth and back.
\begin{figure}[htbp]
\centering
\begin{tikzpicture}
\begin{axis}[
  width=12cm, height=6.5cm,
  xlabel={half-power beamwidth $\theta$ (deg, symmetric)},
  ylabel={directivity (dBi)},
  xmode=log,
  xmin=1, xmax=180, ymin=0, ymax=48,
  axis lines=left,
  legend pos=north east,
  legend cell align=left,
  domain=1:180, samples=300,
]
  % D0[dBi] = 10 log10( 41000 / theta^2 ), theta in degrees, symmetric beam.
  \addplot[engblue, very thick] { 10*log10( 41000/(x*x) ) };
  \addlegendentry{$D_0 \approx 41000/\theta^2$}
  % A conservative estimate constant 26000 (accounts for taper/sidelobes).
  \addplot[engred, thick, dashed] { 10*log10( 26000/(x*x) ) };
  \addlegendentry{tapered ($\approx 26000/\theta^2$)}
\end{axis}
\end{tikzpicture}
\caption[Directivity against half-power beamwidth]{Directivity of a
pencil-beam antenna against its half-power beamwidth, taken symmetric in
the two principal planes. A narrower beam packs the radiated power into a
smaller solid angle, so directivity rises as the inverse square of the
beamwidth: halving the beamwidth quadruples the directivity, a gain of
six decibels. The upper curve is the lossless-uniform estimate and the
lower one carries a realistic amplitude taper and sidelobe budget, the
two bracketing any real antenna. The plot lets an engineer read a gain
from a measured beamwidth without integrating a pattern.}
\label{fig:vol04ch10:directivity-beamwidth}
\end{figure}
